% TODO make hw stylesheet
\documentclass[10pt]{article}
\usepackage[letterpaper,top=1in]{geometry}
\usepackage{quiver}
\usepackage[dvipsnames,svgnames]{xcolor}
\usepackage[shortlabels]{enumitem}
\usepackage[framemethod=TikZ]{mdframed}
\usepackage{amsmath,amssymb,amsthm}
\usepackage[colorlinks]{hyperref}
\usepackage{multicol}
\usepackage{tabularx}

\hypersetup{
	colorlinks=true,
	linkcolor=blue,
	filecolor=magenta,
	urlcolor=magenta,
	pdftitle={MAT 534 HW}
}

\usepackage{mathtools}
\usepackage{thmtools}
\usepackage[textsize=footnotesize]{todonotes}
\usepackage{titling}

\reversemarginpar
\mdfdefinestyle{mdbluebox}{roundcorner=10pt,innerbottommargin=9pt,
	linecolor=blue,backgroundcolor=TealBlue!5,}
\declaretheoremstyle[headfont=\sffamily\bfseries\color{MidnightBlue},
mdframed={style=mdbluebox},]{thmbluebox}

\mdfdefinestyle{mdbloobox}{frametitlefont=\bfseries,innerbottommargin=8pt,
	nobreak=true,backgroundcolor=TealBlue!5,linecolor=blue,}
\declaretheoremstyle[headfont=\bfseries\color{MidnightBlue},
mdframed={style=mdbloobox},headpunct={\\[3pt]},postheadspace=0pt,]{thmbloobox}

\mdfdefinestyle{mdredbox}{frametitlefont=\bfseries,innerbottommargin=8pt,
	nobreak=true,backgroundcolor=Salmon!5,linecolor=RawSienna,}
\declaretheoremstyle[headfont=\bfseries\color{RawSienna},
mdframed={style=mdredbox},headpunct={\\[3pt]},postheadspace=0pt,]{thmredbox}

\mdfdefinestyle{mdgreenbox}{linecolor=ForestGreen,backgroundcolor=ForestGreen!5,
	linewidth=2pt,rightline=false,leftline=true,topline=false,bottomline=false,}
\declaretheoremstyle[headfont=\bfseries\sffamily\color{ForestGreen!70!black},
mdframed={style=mdgreenbox},headpunct={ --- },]{thmgreenbox}

\mdfdefinestyle{mdorangebox}{linecolor=RedOrange,backgroundcolor=RedOrange!5,
	linewidth=2pt,rightline=false,leftline=true,topline=false,bottomline=false,}
\declaretheoremstyle[headfont=\bfseries\sffamily\color{RedOrange!70!black},
mdframed={style=mdorangebox},headpunct={ --- },]{thmorangebox}

\mdfdefinestyle{mdblackbox}{linecolor=black,backgroundcolor=RedViolet!5!gray!5,
	linewidth=3pt,rightline=false,leftline=true,topline=false,bottomline=false,}
\declaretheoremstyle[mdframed={style=mdblackbox}]{thmblackbox}

\declaretheoremstyle[spaceabove=6pt,spacebelow=6pt]{basehead}

\declaretheorem[style=basehead,name=Problem]{problem}
\declaretheorem[style=thmredbox,name=Problem,numbered=no]{problem*}
\declaretheorem[style=thmbloobox,name=Setup]{setup} % for config geo
\declaretheorem[style=thmredbox,name=Example,sibling=problem]{example}
\declaretheorem[style=thmredbox,name=$\bigstar$ Problem,sibling=problem]{niceprob}
\declaretheorem[style=thmbluebox,name=Theorem,sibling=problem]{theorem}
\declaretheorem[style=thmbluebox,name=Corollary,sibling=theorem]{corollary}
\declaretheorem[style=thmbluebox,name=Proposition,sibling=theorem]{prop}
\declaretheorem[style=thmbluebox,name=Lemma,numberwithin=problem]{lemma}
\declaretheorem[style=thmbluebox,name=Theorem,numbered=no]{theorem*}
\declaretheorem[style=thmbluebox,name=Lemma,numbered=no]{lemma*}
\declaretheorem[style=thmgreenbox,name=Claim,numberwithin=problem]{claim}
\declaretheorem[style=thmgreenbox,name=Claim,numbered=no]{claim*}
\declaretheorem[style=thmblackbox,name=Remark,numberwithin=problem]{remark}
\declaretheorem[style=thmblackbox,name=Remark,numbered=no]{remark*}
\declaretheorem[style=thmgreenbox,name=Definition,sibling=theorem]{definition}
\declaretheorem[style=thmgreenbox,name=Definition,numbered=no]{definition*}

\providecommand{\ol}{\overline}
\providecommand{\quiv}{\equiv}
\providecommand{\ceil}[1]{\left\lceil #1 \right\rceil}
\providecommand{\floor}[1]{\left\lfloor #1 \right\rfloor}
\newcommand{\dd}{\mathrm{d}}
\providecommand{\ul}{\underline}
\providecommand{\eps}{\varepsilon}
\providecommand{\half}{\frac{1}{2}}
\providecommand{\CC}{\mathbb C}
\providecommand{\FF}{\mathbb F}
\providecommand{\NN}{\mathbb N}
\providecommand{\QQ}{\mathbb Q}
\providecommand{\RR}{\mathbb R}
\providecommand{\ZZ}{\mathbb Z}
\providecommand{\dg}{^\circ}
\providecommand{\ii}{\item}
\providecommand{\setdiff}{\smallsetminus}
\providecommand{\alert}[1]{{\sffamily\textbf{\textcolor{blue}{#1}}}}
\providecommand{\inv}{^{-1}}
\providecommand{\mb}[1]{\mathbf{#1}}

\newenvironment{soln}{\begin{proof}[Solution]}{\end{proof}}

\providecommand{\pow}{\operatorname{Pow}}
\providecommand{\aut}{\operatorname{Aut}}
\providecommand{\vocab}[1]{{\textbf{\textcolor{ForestGreen}{#1}}}}
\providecommand{\lcm}{\operatorname{lcm}}
\providecommand{\abs}[1]{\left|#1\right|}
\providecommand{\Ob}{\operatorname{Ob}}
\providecommand{\Hom}{\operatorname{Hom}}
\providecommand{\Aut}{\operatorname{Aut}}
\providecommand{\id}{\operatorname{id}}
\newcommand{\doubleparen}[1]{(\!(#1)\!)}

\usepackage{mathrsfs}

% place title at top right
\setlength{\droptitle}{-8ex}
\pretitle{\begin{flushright}\Large\bfseries}
\posttitle{\par\end{flushright}}
\preauthor{\begin{flushright}}
\postauthor{\end{flushright}}
\predate{\begin{flushright}}
\postdate{\end{flushright}}

\title{MAT 534 HW05}
\author{\large Aditya Pahuja}
\date{\large Due 10/08}
\begin{document}
\maketitle

\begin{problem}
    Let $\mathscr C$ be a category.
    \begin{enumerate}[(a)]
	\ii Show that for any object $X\in\operatorname{Ob}(\mathscr C)$,
	there is only one morphism from $X$ to itself
	that has the defining property of $\id_X$.
	\ii Suppose that $f\in\Hom(X,Y)$ is an isomorphism.
	Does this imply that $f$ has a unique inverse in $\Hom(Y,X)$?
	\ii Show that the group $\Aut(X)$
	acts on the set of isomorphisms in $\Hom(X,Y)$
	and that this action is transitive.
    \end{enumerate}
\end{problem}

\begin{soln}
    (a) Let $\id_X$ and $\id'_X$ be two such morphisms. Then
    \begin{equation*}
        \id_X = \id_X\circ\id'_X = \id'_X.
    \end{equation*}

    (b) Yes; if $g,g'\in\Hom(Y,X)$ are two inverses of $f$,
    meaning $g\circ f = g'\circ f = \id_X$ and $f\circ g = f\circ g' = \id_Y$,
    then
    \begin{equation*}
	g = g\circ \underbrace{(f\circ g')}_{\id_Y}
	= \underbrace{(g\circ f)}_{\id_X}\circ g' = g'.
    \end{equation*}

    (c) Define an action of $\Aut(X)$ on the isomorphisms in $\Hom(X,Y)$ by
    $f\cdot g = f\circ g$. This is a right action because
    \begin{equation*}
	f\cdot(g\circ h) = f\circ(g\circ h) = (f\circ g)\circ h
	= (f\cdot g)\cdot h,\quad
	f\cdot\id_X = f\circ\id_X = f.
    \end{equation*}
    The action is transitive because
    given $f_1,f_2\in\Hom(X,Y)$ are isomorphisms then
    $f_2\inv \circ f_1\in\Aut(X)$ satisfies
    \begin{equation*}
        f_2\cdot(f_2\inv\circ f_1) = f_2\circ f_2\inv\circ f_1 = f_1.\qedhere
    \end{equation*}
\end{soln}

\begin{problem}
    Given morphism $f\colon X\to Y$, denote by $(*)$ the property that
    if $g\colon Y\to Z$ and $h\colon Y\to Z$ are two morphisms
    such that $g\circ f = h\circ f$, then $g = h$.
    \begin{enumerate}[(a)]
	\ii Show that in \doubleparen{sets}, $f\colon X\to Y$ is surjective
	iff $f$ satisfies $(*)$.
	\ii Characterize the morphisms satisfying $(*)$ in
	\doubleparen{topological spaces}.
	\ii Characterize the morphisms satisfying $(*)$ in
	\doubleparen{groups}.
	\ii Characterize the morphisms satisfying $(*)$ in
	\doubleparen{abelian groups}.
    \end{enumerate}
\end{problem}

\begin{soln}
    (a) If $f$ is surjective, then for each $y\in Y$,
    there exists $x\in X$ such that $y = f(x)$, so
    $g(y) = g(f(x)) = h(f(x)) = h(y)$ for each $y\in Y$,
    which means $g = h$.

    On the other hand if $f$ isn't surjective, then
    there exists $y_0\in Y$ that is not mapped to by any $x\in X$.
    If $g = \id_Y$ and $h$ is the function mapping $y_0$
    to any other point in $Y$ while fixing all the other points
    (i.e. $h(y_0) = y_1$ for some $y_1\ne y_0$ and
    $h|_{Y\setminus\{y_0\}} = \id_{Y\setminus\{y_0\}}$).
    Then, $g\ne h$ but $g\circ f = h\circ f$
    because $y_0$ is never passed as a value to either $g$ or $h$.
    Thus if $f$ isn't surjective then $(*)$ doesn't hold.
    \\

    (b) \todo{todo}
\end{soln}

\begin{problem}
    Do the following categories have an inital or terminal object?
    If so, describe them; if not, explain why not.
    \begin{enumerate}[(a)]
	\ii \doubleparen{sets}
	\ii \doubleparen{groups}
	\ii \doubleparen{fields}
	\ii \doubleparen{fields of characteristic $0$}
    \end{enumerate}
\end{problem}

\begin{soln}
    (a) \doubleparen{sets} has initial object $\varnothing$,
    since there is tautologically one function $\varnothing\to A$
    for any set $A$ (and no other set is an initial object
    because a morphism $X\to \{1,2\}$ necessarily has
    at least two possible functions by sending everything
    to either $1$ or $2$). Any singleton set is a final object
    because a function from $X$ to a singleton set
    leaves exactly one choice for each element of $X$
    (and no other set is a final object because
    a morphism $f\colon\{1\}\to A$ either has no choices for $f(1)$ if $A$ is empty,
    or at least two choices for $f(1)$ if $A$ has more than one element).
    \\

    (b) Any group with one element is both an initial and final object,
    since a morphism of groups always maps the identity of one
    to the identity of the other. Since all initial objects are isomorphic
    and all final objects are isomorphic (by Problem 6),
    this characterizes all initial and final objects in \doubleparen{groups}.
    \\

    (c) The category of fields does not have an initial object
    or a final object: any morphism of fields is injective,
    so there is no morphism between fields of different characteristics,
    i.e. given any field $F$ there exists another field $F'$
    (for example any field with a different characteristic)
    for which $\Hom(F,F')$ is empty, which automatically means
    $\Hom(F,G)$ cannot have exactly one element for all fields $G$.
    \\

    (d) The category of fields with characteristic zero has
    (any field isomorphic to) $\QQ$ as an initial object,
    since a morphism $\QQ\to F$ is uniquely determined due to the fact
    that $1$ must map to $1_F$.
    On the other hand, there is no final object,
    because apparently you can construct fields of arbitrarily large cardinality
    (I don't understand why well enough to write down why),
    and if $F,F'$ are fields such that
    $F'$ has larger cardinality than $F$ then there is no morphism $F'\to F$
    because field homomorphisms must be injective,
    so $F$ cannot be final.
\end{soln}

\begin{problem}
    In some categories, the objects are sets and have elements,
    and then we can try to describe those elements category-theoretically.
    For example:
    \begin{enumerate}[(a)]
	\ii Let $\mathbf 1$ be any one-element set in \doubleparen{sets},
	for example $\{\varnothing\}$.
	Show that the elements of a set $X$ are in 1-to-1 correspondence
	with the morphisms $\Hom(\mathbf 1, X)$.
	\ii In \doubleparen{groups}, show that
	the elements of a group $G$ are in 1-to-1 correspondence
	with $\Hom(\ZZ,G)$.
	\ii Let $F$ be a field, and consider the category
	\doubleparen{$F$-vector spaces} of vector spaces
	over the field $F$. Can you find a vector space $V_0$ in this category
	such that for any vector space $V$ over $F$,
	the elements of $V$ are in 1-to-1 correspondence with $\Hom(V_0,V)$?
    \end{enumerate}
\end{problem}

\begin{soln}
    (a) Since $\mb 1 = \{\varnothing\}$ has only one element, a morphism
    $f\colon\mb 1\to X$ is uniquely determined by $f(\varnothing)$.
    Thus there is a correspondence between elements $x\in X$
    and morphisms $f\in\Hom(\mb 1,X)$ given by
    \begin{equation*}
        x \longleftrightarrow (\varnothing\mapsto x).
    \end{equation*}

    (b) Since $\ZZ$ is generated by $1$, a morphism
    $f\colon\ZZ\to G$ is uniquely determined by $f(1)$,
    as this tells us $\phi(n) = g^n$ for all $n\in\ZZ$.
    Thus there is a correspondence between elements $g\in G$
    and morphisms $\phi\in\Hom(\ZZ,G)$ given by
    \begin{equation*}
	g\longleftrightarrow (n\mapsto g^n).
    \end{equation*}

    (c) Let $V_0 = F$. Since $V_0$ is a $1$-dimensional $F$-vector space,
    a morphism $f\in\Hom(V_0,V)$ is determined uniquely by $f(1)$.
    Thus there is a correspondence between elements $v\in V$
    and morphisms $f\in\Hom(V_0,V)$ given by
    \begin{equation*}
        v\longleftrightarrow (c\mapsto cv).\qedhere
    \end{equation*}
\end{soln}

\begin{problem}
    Let $\mathscr C$ be a category, and $A\in\mathscr C$ a fixed object.
    For every $X\in\Ob(\mathscr C)$, let $F(X) = \Hom(A,X)$
    be the set of morphisms from $A$ to $X$ in our category.
    Show that $F$ is a functor from $\mathscr C$ to \doubleparen{sets}.
\end{problem}

\begin{soln}
    If $f\colon X\to Y$ is a morphism in $\mathscr C$,
    let $F(f)$ be the morphism $F(X)\to F(Y)$ sending
    \begin{equation*}
	(g\colon A\to X)\mapsto(f\circ g\colon A\to Y)
    \end{equation*}
    where $g\in\Hom(A,X)$ is any morphism. This is a functor because
    \begin{itemize}
        \ii $F(\id_X)$ is the map
	$(g\colon A\to X)\mapsto(\id_X\circ g = g\colon A\to X)$
	which is $\id_{F(X)}$.
	\ii If $f_1\in\Hom(X,Y)$ and $f_2\in\Hom(Y,Z)$ then
	$F(f_2)\circ F(f_1)$ is the composed map
	\begin{equation*}
	    (g\colon A\to X)\mapsto(f_1\circ g\colon A\to Y)
	    \mapsto(f_2\circ f_1\circ g\colon A\to Z)
	\end{equation*}
	where $g\in\Hom(A,X)$, whose net effect is
	exactly that of $F(f_2\circ f_1)\colon F(X)\to F(Z)$,
	i.e. $F$ is well-behaved under composition.\qedhere
    \end{itemize}
\end{soln}

\begin{problem}
    Let $A$ and $B$ be two initial (or final) objects in a category.
    Show that there is a unique isomorphism between $A$ and $B$.
\end{problem}

\begin{soln}
    If $A$ and $B$ are two initial objects
    or two final objects, then the important data is that
    $\Hom(A,A)$, $\Hom(A,B)$, $\Hom(B,A)$, and $\Hom(B,B)$ each have
    exactly one morphism, which we will respectively call
    $f_{AA} = \id_A$, $f_{AB}$, $f_{BA}$, and $f_{BB} = \id_B$. Then
    \begin{align*}
	f_{AB}\circ f_{BA}\in\Hom(B,B)&\implies f_{AB}\circ f_{BA} = \id_B,\\
	f_{BA}\circ f_{AB}\in\Hom(A,A)&\implies f_{BA}\circ f_{AB} = \id_A,
    \end{align*}
    so $f_{AB}$ is an isomorphism,
    and in fact it is the unique (iso)morphism from $A$ to $B$.
\end{soln}

\begin{problem}
    Given a topological space $X$, we can define
    a category $\mathscr C_X$ as follows:
    The objects of $\mathscr C_X$ are the open subsets of $X$,
    and for any two open subsets $U,V\subseteq X$,
    the set of morphisms in $\mathscr C_X$ is
    \begin{equation*}
        \Hom(U,V) = \begin{cases}
	    \{i_{UV}\} & \text{if }U\subseteq V,\\
	    \varnothing & \text{if }U\not\subseteq V.
        \end{cases}
    \end{equation*}
    If $U\subseteq V\subseteq W$, we define composition as
    $i_{VW}\circ i_{UV} = i_{UW}$.
    \begin{enumerate}[(a)]
        \ii Prove that $\mathscr C_X$ is indeed a category.
	\ii Does $\mathscr C_X$ have an initial or final object?
	\ii For every open $U\subseteq X$, define $F(U)$ to be
	the set of continuous functions from $U$ to $\RR$.
	Show that $F$ is a contravariant functor from $\mathscr C_X$
	to the category \doubleparen{$\RR$-vector spaces}.
    \end{enumerate}
\end{problem}

\begin{soln}
    (a) To show that $\mathscr C$ is a category, we only need to check
    that the morphisms behave as they are supposed to.
    Each object has identity $i_{UU}$
    (because inclusion of $U$ into $U$ is the identity), and
    if $U\subseteq V\subseteq W\subseteq Z$ are open subsets then
    \begin{equation*}
	(i_{WZ}\circ i_{VW})\circ i_{UV} =
	i_{VZ}\circ i_{UV} =
	i_{UZ} =
	i_{WZ}\circ i_{UW} =
	i_{WZ}\circ(i_{VW}\circ i_{UV}),
    \end{equation*}
    proving associativity.
    \\

    (b) Each $\Hom(U,V)$ has at most one morphism, so
    the only way that $U$ is initial is if $U\subseteq V$ for all open $V$,
    i.e. $U = \varnothing$, so $\boxed{\varnothing}$ is the unique initial object;
    and the only way that $V$ is final is if $V\supseteq U$ for all open $U$,
    i.e. $V = \boxed X$ is the unique final object.
    \\

    (c) If $U,V\subseteq X$ are open sets
    such that $U\subseteq V$, then the map $F(V)\to F(U)$ given by
    \begin{align*}
	F(i_{UV}) : F(V)&\to F(U)\\
	(f\colon V\to\RR)&\mapsto (f\circ i_{UV}\colon U\to\RR)
    \end{align*}
    is well-defined because inclusion is a continuous map
    and is obviously linear since $(f+g)\circ i_{UV} = f\circ i_{UV} + g\circ i_{UV}$
    and $(cf)\circ i_{UV} = c(f\circ i_{UV})$ essentially by definition,
    so we will say $F$ sends $i_{UV}$
    to this morphism between $\RR$-vector spaces.

    Now, we check that $F$ is a contravariant functor:
    \begin{itemize}
	\ii For any open $U\subseteq X$, $F(\id_U) = F(i_{UU})$
	is the map $(f\colon U\to\RR)\mapsto(f\circ i_{UU} = f\colon U\to\RR)$,
	so it is the identity on $F(U)$.
	\ii If $U\subseteq V\subseteq W$ are open sets in $X$, then
	$F(i_{UV})\circ F(i_{VW})$ is the composed map
	\begin{equation*}
	    (f\colon W\to\RR)\mapsto(f\circ i_{VW}\colon V\to\RR)
	    \mapsto(f\circ i_{VW}\circ i_{UV}\colon U\to\RR)
	\end{equation*}
	which is the same as the map $F(i_{VW}\circ i_{UV})$ by definition.
    \end{itemize}
\end{soln}










\end{document}
